\section{方程式と不等式}\label{節:微分方程式不等式}
方程式の根を式で表せなくても,根が存在することや,根の個数は分かることがある.ここでは,連続性が存在を保証し,増減が個数を制限する.不等式では,左右の差を一つの関数にして,その最小値を調べる.

\begin{example}{方程式と不等式}
\begin{enumerate}[label=(\arabic*),leftmargin=*]
 \item 方程式 $x=\cos x$ が区間 $[0,1]$ にただ一つの解をもつことを示せ.
 \item $x>0$ のとき $\log x\le x-1$ を示し,等号成立条件を求めよ.
\end{enumerate}
\end{example}
(1) では $h(x)=x-\cos x$ とおく.\,$h$ は $[0,1]$ で連続で,
\[
 h(0)=-1<0,\qquad h(1)=1-\cos1>0
\]
である.中間値の定理から,\,$h(x)=0$ となる点が $(0,1)$ に存在する.一方,この区間では $h'(x)=1+\sin x>0$ だから,\,$h$ は狭義に増加する.二つの異なる点でともに $0$ をとることはできず,解はただ一つである.グラフの交わりを予想しただけではなく,存在と一意性を別々に示した.

(2) は右辺から左辺を引き,\,$g(x)=x-1-\log x$ とおく.
\[
 g'(x)=1-\frac1x=\frac{x-1}{x}.
\]
定義域では $x>0$ なので,符号は分子 $x-1$ で決まる.\,$g$ は $0<x<1$ で減少し,\,$x>1$ で増加する.したがって,最小値は $g(1)=0$ であり,\,$g(x)\ge0$ を得る.両側で狭義に増減しているため,等号は $x=1$ のときに限る.


